\documentclass[12pt, twoside]{article}
\input{../preamble}

\fancyhead[LE]{\thepage}
\fancyhead[RO]{Markscheme}
\fancyhead[LO]{La Scuola d'Italia / Dr.\ Huson / IB Math \\* 29 April 2026}

\begin{document}
\subsubsection*{7.4 Classwork: Solving Exponential Equations --- Markscheme}

\begin{enumerate}[itemsep=0.3cm]

%--- Problem 1: Technique drill --- 4 sub-parts ---
\item[1.] [Maximum mark: 8]

\textbf{(a)} $3^x = 17$. \;Take $\log$ of both sides:
$x \log 3 = \log 17$. \hfill \textbf{M1}

$x = \dfrac{\log 17}{\log 3} = 2.5789\ldots = 2.58$ (3 s.f.) \hfill \textbf{A1}

\medskip

\textbf{(b)} $4^{\,x-2} = 30$. \;$(x - 2) \log 4 = \log 30$. \hfill \textbf{M1}

$x - 2 = \dfrac{\log 30}{\log 4} = 2.4534\ldots$, \;so\;
$x = 4.4534\ldots = 4.45$ (3 s.f.) \hfill \textbf{A1}

\medskip

\textbf{(c)} $2^{\,3x+1} = 100$. \;$(3x + 1) \log 2 = \log 100$.
\hfill \textbf{M1}

$3x + 1 = \dfrac{\log 100}{\log 2} = \dfrac{2}{\log 2} = 6.6438\ldots$, \;so\;
$x = \dfrac{6.6438\ldots - 1}{3} = 1.8813\ldots = 1.88$ (3 s.f.)
\hfill \textbf{A1}

\medskip

\textbf{(d)} $5^x = 7^{\,x-1}$. \;Take $\log$:
$x \log 5 = (x - 1)\log 7$. Collect $x$-terms:
$x(\log 7 - \log 5) = \log 7$. \hfill \textbf{M1}

$x = \dfrac{\log 7}{\log 7 - \log 5}
= \dfrac{0.8451}{0.8451 - 0.6990}
= 5.7853\ldots = 5.79$ (3 s.f.) \hfill \textbf{A1}

\newpage

%--- Problem 2: Bacteria growth ---
\item[2.] [Maximum mark: 7]

$P(t) = 200 \cdot 3^{\,t/10}$.

\textbf{(a)} At $t = 0$: $P(0) = 200 \cdot 3^0 = 200$. \hfill \textbf{A1}

\medskip

\textbf{(b)} $P(24) = 200 \cdot 3^{\,24/10} = 200 \cdot 3^{2.4}$. \hfill \textbf{M1}

$= 200 \cdot 13.9666\ldots = 2793.3\ldots \approx 2793$ \hfill \textbf{A1}

\medskip

\textbf{(c)} $200 \cdot 3^{\,t/10} = 5000 \;\Rightarrow\; 3^{\,t/10} = 25$
\hfill \textbf{M1\;A1}

Take $\log$: \;$\dfrac{t}{10} \log 3 = \log 25$, \;so\;
$\dfrac{t}{10} = \dfrac{\log 25}{\log 3} = 2.9299\ldots$ \hfill \textbf{M1}

$t = 29.299\ldots \approx 29.3$ hours (1 d.p.) \hfill \textbf{A1}

\newpage

%--- Problem 3: Radioactive decay (iodine-131) ---
\item[3.] [Maximum mark: 7]

$M(t) = 50 \cdot \left(\dfrac{1}{2}\right)^{t/8}$.

\textbf{(a)} $M(20) = 50 \cdot (0.5)^{20/8} = 50 \cdot (0.5)^{2.5}$.
\hfill \textbf{M1}

$= 50 \cdot 0.17677\ldots = 8.8388\ldots \approx 8.84$ g (3 s.f.)
\hfill \textbf{A1}

\medskip

\textbf{(b)} $60\,\%$ decayed means $40\,\%$ remains, so $M(t) = 20$:

$50 \cdot (0.5)^{t/8} = 20 \;\Rightarrow\; (0.5)^{t/8} = 0.40$
\hfill \textbf{M1}

Take $\log$: \;$\dfrac{t}{8} = \dfrac{\log 0.40}{\log 0.5}
= \dfrac{-0.39794}{-0.30103} = 1.3219\ldots$ \hfill \textbf{M1}

$t = 8 \cdot 1.3219\ldots = 10.575\ldots \approx 10.6$ days (3 s.f.)
\hfill \textbf{A1}

\medskip

\textbf{(c)} $50 \cdot (0.5)^{t/8} = 1 \;\Rightarrow\; (0.5)^{t/8}
= \dfrac{1}{50} = 0.02$ \hfill \textbf{M1}

$\dfrac{t}{8} = \dfrac{\log 0.02}{\log 0.5} = 5.6439\ldots$, \;so\;
$t = 45.15\ldots \approx 45$ days \hfill \textbf{A1}

\newpage

%--- Problem 4: Cooling ---
\item[4.] [Maximum mark: 6]

$T(t) = 20 + 65 \cdot (0.9)^t$.

\textbf{(a)} $T(0) = 20 + 65 \cdot 1 = 85\,^{\circ}\text{C}$.
\hfill \textbf{A1}

\medskip

\textbf{(b)} $T(15) = 20 + 65 \cdot (0.9)^{15}$. \hfill \textbf{M1}

$= 20 + 65 \cdot 0.20589\ldots = 20 + 13.383\ldots
= 33.383\ldots \approx 33.4\,^{\circ}\text{C}$ (3 s.f.) \hfill \textbf{A1}

\medskip

\textbf{(c)} Set $T(t) = 40$:

$20 + 65 \cdot (0.9)^t = 40 \;\Rightarrow\; (0.9)^t
= \dfrac{20}{65} = \dfrac{4}{13} \approx 0.30769$ \hfill \textbf{M1}

Take $\log$: \;$t \log 0.9 = \log\!\left(\dfrac{4}{13}\right)$, \;so\;
$t = \dfrac{\log(4/13)}{\log 0.9} = \dfrac{-0.51217}{-0.04576}$
\hfill \textbf{M1}

$t = 11.193\ldots \approx 11.2$ minutes (3 s.f.) \hfill \textbf{A1}

\end{enumerate}
\end{document}
